\documentclass[11pt]{article}
\usepackage[margin=1in]{geometry}
\usepackage{amsmath,amssymb,amsthm}
\usepackage[utf8]{inputenc}
\usepackage{listings}
\usepackage{hyperref}
\usepackage{xcolor}
\usepackage{textcomp}

\definecolor{leanblue}{rgb}{0.2,0.4,0.7}

\lstdefinelanguage{Lean}{
  keywords={theorem,lemma,def,axiom,structure,class,instance,where,by,sorry,simp,exact,intro,apply,have,let,fun,match,if,then,else,import,open,namespace,end,section,variable,inductive,noncomputable,abbrev,deriving},
  keywordstyle=\color{leanblue}\bfseries,
  comment=[l]{--},
  morecomment=[s]{/-}{-/},
  string=[b]",
  basicstyle=\ttfamily\small,
  breaklines=true,
  extendedchars=true,
  inputencoding=utf8,
  literate=
    {≥}{{$\geq$}}1
    {≤}{{$\leq$}}1
    {→}{{$\to$}}1
    {←}{{$\leftarrow$}}1
    {∧}{{$\wedge$}}1
    {∨}{{$\vee$}}1
    {¬}{{$\neg$}}1
    {∀}{{$\forall$}}1
    {∃}{{$\exists$}}1
    {⊆}{{$\subseteq$}}1
    {⊂}{{$\subset$}}1
    {∈}{{$\in$}}1
    {∉}{{$\notin$}}1
    {×}{{$\times$}}1
    {⊗}{{$\otimes$}}1
    {▸}{{$\blacktriangleright$}}1
    {⟨}{{$\langle$}}1
    {⟩}{{$\rangle$}}1
    {λ}{{$\lambda$}}1
    {α}{{$\alpha$}}1
    {β}{{$\beta$}}1
    {σ}{{$\sigma$}}1
    {θ}{{$\theta$}}1
    {ε}{{$\varepsilon$}}1
    {μ}{{$\mu$}}1
    {π}{{$\pi$}}1
    {Σ}{{$\Sigma$}}1
    {Π}{{$\Pi$}}1
    {▶}{{$\triangleright$}}1
    {⊔}{{$\sqcup$}}1
    {⊓}{{$\sqcap$}}1
    {⊥}{{$\bot$}}1
    {⊤}{{$\top$}}1
    {∅}{{$\emptyset$}}1
    {∪}{{$\cup$}}1
    {∩}{{$\cap$}}1
    {≠}{{$\neq$}}1
    {ℕ}{{$\mathbb{N}$}}1
    {₀}{{$_0$}}1
    {₁}{{$_1$}}1
    {₂}{{$_2$}}1
    {|>}{{$|{\!>}$}}2
    {·}{{$\cdot$}}1
    {—}{{---}}1
    {∘}{{$\circ$}}1
    {√}{{$\sqrt{}$}}1
    {∝}{{$\propto$}}1
    {≈}{{$\approx$}}1
    {═}{{=}}1
    {⟹}{{$\Longrightarrow$}}1,
}

\newtheorem{theorem}{Theorem}
\newtheorem{lemma}[theorem]{Lemma}
\newtheorem{definition}[theorem]{Definition}
\newtheorem{axiom_stmt}[theorem]{Axiom}

\title{Formal Verification: Prism Peer Sampling via Fisher-Yates}
\author{Zach Kelling, Lux Industries}
\date{December 2025}

\begin{document}
\maketitle

\begin{abstract}
We formalize Prism, the peer sampling protocol that implements Fisher-Yates shuffle with cryptographic randomness. We state the uniformity of Fisher-Yates as an axiom (referencing Knuth Vol.~2, Algorithm~P) and prove that selecting the first $k$ elements gives a uniform $k$-subset.
\end{abstract}

\section{Introduction}
Prism provides the Cut interface for peer sampling. Given a population of peers, it samples $k$ uniformly at random using \texttt{crypto/rand}. The Fisher-Yates uniformity is a well-known result that we axiomatize.

\section{Definitions}
\begin{definition}[Luminance]
Network health metric: active peers, total peers, and illuminance ratio.
\end{definition}


\section{Theorems and Axioms}
\begin{axiom_stmt}[\texttt{fisher\_yates\_uniform}]
Fisher-Yates shuffle produces each of $n!$ permutations with equal probability (Knuth Vol.~2, Algorithm~P).
\end{axiom_stmt}

\begin{axiom_stmt}[\texttt{prefix\_k\_is\_uniform\_subset}]
Selecting the first $k$ elements from a uniform permutation gives a uniform $k$-subset: $\binom{n}{k} > 0$.
\end{axiom_stmt}

\begin{theorem}[\texttt{sufficient\_luminance}]
If $k \leq \text{activePeers}$ and $\text{activePeers} \leq \text{totalPeers}$, then $k \leq \text{totalPeers}$.
\end{theorem}
\begin{proof}
By \texttt{omega}.
\end{proof}

\begin{theorem}[\texttt{more\_active\_more\_committees}]
More active peers means more possible committees: $\binom{n}{k} \geq 1$ when $k \leq n$.
\end{theorem}


\section{Lean Source}
\lstinputlisting[language=Lean]{../../formal/lean/Protocol/Prism.lean}

\section{References}
\noindent Source code: \url{https://github.com/luxfi/proofs/blob/main/lean/Protocol/Prism.lean}\\
\noindent White papers: \url{https://papers.lux.network}

\noindent Related white papers:
\begin{itemize}
  \item \texttt{lux-consensus.tex}
\end{itemize}

\end{document}
